\documentclass[10pt,a4paper]{amsart}
\usepackage[margin=19mm]{geometry}
\usepackage{amsmath,amssymb,mathtools}
\usepackage[scr=rsfs,cal=euler]{mathalfa}
\usepackage{xfrac}
\usepackage[unicode,hidelinks,bookmarks=false]{hyperref}
\newcommand{\ie}{i.e.\ }
\newcommand{\cf}{cf.\ }
\newcommand{\resp}{respectively\ }
\newcommand{\Subsection}{\subsection}
\providecommand{\nobreakdash}{\nobreak\hbox{-}}
\setlength{\emergencystretch}{2em}
\multlinegap0pt
\usepackage{mathtools}
\usepackage{tikz}
\newtheorem{proposition}{Proposition}
\newcommand{\bfM}{{\mathbf{M}}}
\newcommand{\brsq}[1]{{\left[{#1}\right]}}
\newcommand{\defeq}{{\overset{\triangle}{=}}}
\newcommand{\mX}{{\mathcal{X}}}
\newcommand{\obfM}{{\overline{\mathbf{M}}}}
\newcommand{\set}[1]{{\left\{{#1}\right\}}}
\newcommand{\tQ}{{\widetilde{Q}}}
\newcommand{\ve}{\varepsilon}
\newcommand{\vp}{\varphi}
\title{Simple Finite-Length Achievability and Converse Bounds for Deletion and Insertion Channels}
\author{Ruslan Morozov, Tolga Mete Duman}
\date{Source: arXiv:2504.20961v4 (CC BY 4.0)}
\begin{document}
\maketitle

\noindent\emph{Source and licence.} Simple Finite-Length Achievability and Converse Bounds for Deletion and Insertion Channels. Version arXiv:2504.20961v4 (CC BY 4.0), distributed by arXiv under the Creative Commons Attribution 4.0 International licence, http://creativecommons.org/licenses/by/4.0/, as recorded in the arXiv licence metadata for this exact version and shown on the abstract page https://arxiv.org/abs/2504.20961v4. Full licence notice: \texttt{licenses/arxiv-2504.20961-CC-BY-4.0.txt}.\par\smallskip

\noindent\textit{Editorial scope.} Counted unit: the article's proposition p:mocvb (source0.tex lines 714--721). The article's complete original proof of it is delivered with it (source0.tex lines 723--752, one \texttt{proof} environment of 30 lines). No mathematics is written, repaired or re-derived: the statement and the proof are the article's own lines, in the article's own notation, and no retained line is substituted or re-flowed.  Retained uncounted as the source-local context the proof needs: lines 663--668.  Nothing was omitted from any retained span, so the package carries no omission record.  No retained line contains a citation, so the wrapper carries no bibliography and no bibliography entry.  No retained line contains a figure environment, an image-inclusion command or drawing code, so no image is delivered and the package declares no asset dependency.  Editorial formatting only, in addition: an AMS \texttt{amsart} preamble that restates the article's own declarations and macros used by the retained lines, namely the environments \texttt{proposition} on separate counters, together with the standard packages those lines need.  The retained environments print in this document as Proposition 1 in order of appearance, whereas the article prints the same items with the numbering of its own full document.  The complete native source of the article as distributed in the arXiv e-print for this exact version is bundled unchanged as \texttt{sources/177293/source0.tex}.
\begin{align}
&\bfM(W,\ve) \leq \obfM_{\text{SW}}(W,\ve,Q,\vp)
\nonumber \\&
\ \ \ \defeq \frac{1}{\vp - \ve} \cdot \max_{x\in\mX} \sup \set{\rho\bigg| \Pr_{W(Y|x)}\brsq{\frac{W(Y|x)}{Q(Y)}\leq \rho} \leq \vp}
\label{eq:cvb}
\end{align}

\begin{proposition}
\label{p:mocvb}
The maximum code size which allows transmission over channel $W$ with error probability $\ve$ is upper-bounded by
\begin{align}
\bfM(W, \ve) \leq \frac{\tau(W)}{1-\ve}.
\label{eq:p1w}
\end{align}
\end{proposition}

\begin{proof}
Consider the case of $\vp=1$ in \eqref{eq:cvb}.
The $1$-quantile is just the maximum possible value of $W(y|x)/Q(y)$:
\begin{align}
&\obfM_{\text{SW}}(W,\ve,Q,1) 
\nonumber\\ & \ \ =
 \frac{1}{1 - \ve} \cdot \max_x\inf \set{\rho\bigg| \Pr_{W(Y|x)}\brsq{\frac{W(Y|x)}{Q(Y)}\leq  \rho}= 1}
\nonumber\\ &
\ \ = \frac{1}{1 - \ve} \cdot \max_y \frac{\max_x W(y|x)}{Q(y)}.
\label{eq:mvp1}
\end{align}
Due to the constraint $\sum_y Q(y)=1$, the minimum of \eqref{eq:mvp1} over $Q$ is achieved, when $\max_x W(y|x)/Q(y)$ is the same for all probable\footnote{If $W(y|x)=0$, then $y$ does not contribute to $\Pr[..]$ in \eqref{eq:mvp1}.} values of $y$, i.e., when $Q(y)$ is proportional to $\max_x W(y|x)$.
Denote such distribution by $\tQ_W$: 
{\allowdisplaybreaks
\begin{align}
\tQ_W(y)&\;\defeq\;\frac{\max_x W(y|x)}{\tau(W)}.
\label{eq:tQ}
\end{align}
}%
Note that it is optimal for $\vp=1$, i.e., 
\begin{align}
\obfM_\text{SW}(W,\ve,\tQ_W,1)=\inf_Q\obfM_\text{SW}(W,\ve,Q,1).
\label{eq:tQopt}
\end{align}
Substituting distribution $\tQ_W$ in \eqref{eq:mvp1}, one obtains
\begin{align}
\bfM(W,\ve)\leq \obfM_\text{SW}(W,\ve,\tQ_W,1)=\frac{\tau(W)}{1-\ve}.
\label{eq:cvbs}
\end{align}
\end{proof}

\end{document}
